\documentclass[aspectratio=169,english,plain,c]{beamer}

\usepackage{lib/beamer-themes/dividing-lines/beamerthemedividing-lines}

\usepackage[T1]{fontenc}
\usepackage[utf8]{inputenc}
\usepackage{tikz}
\usepackage{fontawesome}
\usepackage{etoolbox}
\usetikzlibrary{backgrounds,graphs,arrows.meta}

\definecolor{BaseGray}{RGB}{66,66,66} % rgb(66,66,66)

\colorlet{SoftGray}{BaseGray!40}
\colorlet{BackGray}{BaseGray!5}
\colorlet{SoftTextGray}{BackGray!60!SoftGray}

\tikzset{
   FunctionDef/.style={
      draw=BaseGray,
      fill=BaseGray,
      minimum width=1.55cm,
      minimum height=1cm,
      text=white,
      font=\bfseries,
      text centered,
      inner sep=0pt,
      rounded corners=1mm,
      outer sep=2pt
   },
   Blob/.style={
      draw=SoftGray,
      fill=SoftGray,
      minimum size=4mm,
      text=white,
      circle,
      font=\bfseries,
      text centered,
      inner sep=0pt,
      outer sep=2pt
   },
   Def/.style={
      Blob,
      rectangle, rounded corners=1mm
   },
   ActiveBlob/.style={
      Blob,
      draw=BaseGray, fill=BaseGray
   },
   FunctionBack/.style={
      fill=BackGray,
      % draw=SoftGray,
      rounded corners=2mm,
      rectangle
   },
   Link/.style={
      draw=SoftGray,
      line width=1.5pt,
      line cap=round,
      line join=round,
      -%
   },
   FuncLink/.style={
      Link,
      draw=SoftGray,
      dotted
   },
   Cursor/.style={
      fill=BackGray,
      draw=BaseGray,
      line join=round,
      line cap=round
   },
   Hover-Over/.style={
      fill=BackGray,
      draw=SoftGray,
      opacity=.5,
      draw opacity=1,
      rounded corners=1mm,
      line join=round,
      line cap=round
   },
   Line-Of-Text/.style={
      fill=#1,
      draw=none,
      rounded corners=1.5pt,
      inner sep=1pt,
      minimum width=1cm,
      minimum height=6.5pt
   },
   Input-Base/.style={
      fill=BackGray,
      draw=SoftGray,
      rounded corners=1mm,
      inner sep=1pt,
      minimum width=2cm,
      minimum height=12pt
   },
   % code sub-styles
   A/.style={Line-Of-Text=SoftTextGray},
   B/.style={},
   C/.style={Line-Of-Text=SoftGray}
}

\usepackage[english]{babel}
\usepackage{microtype}

\usepackage{etoolbox,tabularx,array}
\makeatletter
\appto\input@path{{lib/color-palettes/}{lib/sopra-collection/sopra-listings}{lib/code-animation/}}

\usepackage{lib/color-palettes/color-palettes}
\colorlet{paletteA}{BaseGray}
\colorlet{paletteB}{BaseGray}
\colorlet{paletteC}{BaseGray}
\colorlet{paletteD}{BaseGray}
\usepackage[cpalette,encoding,defaultfont,fakeminted]{lib/sopra-collection/sopra-listings/sopra-listings}

\solLoadLanguage{R}
\solsetmintedstyle{plain}
\usepackage{lib/code-animation/code-animation}
\SetCodeAnimationFadeOutMixin{!10!btdl@color@background}
\CodeAnimationsWithFadeOut

\SolDefineStyles{%
 {keywordA: \color{sol@colors@lst@keywordA}\bfseries},%
 {keywordB: \color{sol@colors@lst@keywordB}\bfseries},%
 {keywordC: \color{sol@colors@lst@keywordC}\bfseries},%
 {keywordD: \color{sol@colors@lst@keywordC}\bfseries},%
}

\usepackage[bare]{lib/tikzpingus/tex/tikzpingus}
\usepackage{forest}

\def\DoTitlepage{}%
\def\DoEndpage{}%

\title{SE-Projekte im WS 23/24 --- flowR}
\subtitle{Static Program Slicing with flowR}
\author{Florian Sihler}
\date{20.07.2023}

\newsavebox\HappyPingu
\savebox\HappyPingu{\tikz{\pingu[eyes wink,right wing wave,:mix-all=30!btdl@color@background]}}
\newsavebox\FlowR
\savebox\FlowR{\usebeamerfont{title}\sbfamily\includegraphics[height=2.5\baselineskip]{data/flowR.pdf}}

\begin{document}

\begin{frame}
\begin{tikzpicture}[overlay,remember picture]
%    \node[below left=3mm] (@flowr) at(current page.north east) {%
%    \includegraphics[width=2.5cm]{data/flowR.pdf}
%     };
   % fake opacity
   % \fill[btdl@color@background,opacity=.5] (@flowr.north east) rectangle (@flowr.south west);
   \node[right,xshift=5mm] (@) at(current page.west) {%
      \usebeamerfont{title}\sbfamily Static Program Slicing with%
   };
   \node[above right,yshift=-.482\ht\FlowR] (@flowr) at(@.base east) {\usebox\FlowR};
   \fill[btdl@color@background,opacity=.3] (@flowr.north east) rectangle (@flowr.south west);

   \node[below right,yshift=1.25mm,font=\itshape\sbfamily,SoftGray] at (@.south west) {%
      in the winter semester 2023/24%
   };
   \node[above left=2.5mm,yshift=-.65mm] (@) at(current page.south east) {\href{https://github.com/EagleoutIce/tikzpingus}{\scalebox{0.2}{\usebox\HappyPingu}}};
   \node[above right=2.5mm,opacity=.55] at(current page.south west) {\href{https://www.uni-ulm.de/in/sp/}{\includegraphics[height=\dimexpr0.13999\ht\HappyPingu+0.13999\dp\HappyPingu]{data/sp.pdf}}};
\end{tikzpicture}
\end{frame}

\def\SupportedTypes{
   \coordinate (@) at([xshift=-1mm,yshift=-4.5mm]current page.north east);
   \foreach \i in {PSE2, PSE1, AP SE} {
      \node[left=1.5mm,fill=BackGray,draw=BackGray,rounded corners=2pt,text=BaseGray] (@) at(@.west) {\tiny\i};
   }
}


\forestset{T/.style={for tree={font=\ttfamily,align=center,l sep=0pt,l sep-=5mm,child anchor=north}}}
\newsavebox\ArchitectureBox
\setbox\ArchitectureBox=\hbox{\tikz[baseline={([yshift=-2mm]current bounding box.center)},scale=1.35]{\draw[rounded corners=.5mm,fill=btdl@color@background,opacity=.8]
(0,0) -| ++(-6mm,8mm) -- ++(4mm,0) |- ++(2mm,-2mm) coordinate (@) -- cycle
([yshift=-1mm]@) -- ++(0,1mm) -- ++(-2mm,2mm) -- ++(-1mm,0);}}
\newsavebox\Parsing
\setbox\Parsing=\hbox{\scalebox{.5}{\begin{tikzpicture}
   \foreach[count=\i from 0] \Code in {
         {1/C,2/A,3/A},
         {.5/B,1/A,2/A,1/A,1/A},
         {.5/B,2/A,3/A},
         {1/A},
         {},
         {1/A,2/A,6/A},
         {1.5/A,1/A,2/A,.5/B,1/A},
      } {
         \def\XShift{5}
         \foreach \CW/\Style in \Code {
            \ifstrequal{\Style}{B}{\def\RandomSuffix{0}}{\def\RandomSuffix{(rand*0.4mm+0.75mm)}}
            \pgfmathsetmacro\w{3*\CW mm+\RandomSuffix}
            \path[\Style] ([yshift=-2.5mm-1cm,xshift=3.25mm+\XShift pt]0,-\i*0.28*10mm+1mm) rectangle ++(\w pt,-1.75mm);

            \pgfmathsetmacro{\XShift}{\XShift+\w pt+1mm}
            \xdef\XShift{\XShift}
         }
      }
\end{tikzpicture}}}
\newsavebox\FirstAst
\setbox\FirstAst=\hbox{\resizebox*!{1.25cm}{\begin{forest}
   for tree={Blob,edge=Link}
   [[[][,phantom]][[[][,phantom]][]]]
\end{forest}}}
\newsavebox\DataFlow
\setbox\DataFlow=\hbox{\resizebox*!{1.25cm}{\color{lightgray!80!BaseGray}\begin{forest}
   for tree={Blob,edge=Link}
   [[[,name=a][,phantom]][,name=d[[,name=b][,phantom]][,name=c]]]
   \draw[Link,very thick,BaseGray,line cap=round] (a) to[bend right=25] (b) (b) to[bend right=25] (c) (c) to[bend right=25] (d);
\end{forest}}}
\newsavebox\Slicing
\setbox\Slicing=\hbox{\resizebox*!{1.25cm}{\color{lightgray!80!BaseGray}\begin{forest}
   for tree={Blob,edge=Link}
   [[[,name=a][,phantom]][,name=d[[,name=b,ActiveBlob,very thick][,phantom]][,name=c,ActiveBlob,very thick]]]
   \draw[Link,very thick,line cap=round,BaseGray!80!gray] (a) to[bend right=25] (b) (b) to[bend right=25] (c) (c) to[bend right=25] (d);
\end{forest}}}
\newsavebox\Reconstruct
\setbox\Reconstruct=\hbox{\scalebox{.5}{\begin{tikzpicture}
   \foreach[count=\i from 0] \Code in {
         {1/C,2/A,3/A},
         {.5/B,2/A,3/A},
         {1/A},
         {},
         {1/A,2/A,4/A},
      } {
         \def\XShift{5}
         \foreach \CW/\Style in \Code {
            \ifstrequal{\Style}{B}{\def\RandomSuffix{0}}{\def\RandomSuffix{(rand*0.4mm+0.75mm)}}
            \pgfmathsetmacro\w{3*\CW mm+\RandomSuffix}
            \path[\Style] ([yshift=-2.5mm-1cm,xshift=3.25mm+\XShift pt]0,-\i*0.28*10mm+1mm) rectangle ++(\w pt,-1.75mm);

            \pgfmathsetmacro{\XShift}{\XShift+\w pt+1mm}
            \xdef\XShift{\XShift}
         }
      }
\end{tikzpicture}}}
\newsavebox\Overview
\begin{lrbox}\Overview
\scalebox{.8}{\begin{tikzpicture}[br/.style={fill=BackGray!80,minimum width=2.5cm,minimum height=1.5cm,signal,signal from=west,signal to=east,signal pointer angle=125},k/.style={below,font=\footnotesize,xshift=-.33mm,color=darkgray},m/.style={above,font=\scriptsize,xshift=-.33mm,color=gray},a/.style={},BaseGray]
   \tikzset{a/.style={opacity=.5}}
      \scope[transparency group,a]
      \node[br] (@) at (0,0) {};
      \node (file) at (@) {\rotatebox{-5}{\usebox\ArchitectureBox}\llap{\T{.r\kern8pt}}};
      \node[k] at (@.south) {Input File};
      \endscope

      \node[br,right=3mm] (@) at (@.east) {};
      \node (r-conv) at (@) {\usebox\Parsing};
      \node[k] at (@.south) {Parsing};
      % \onslide<4->{\node[m] at (@.north) {\T{parse}\,\(\to\)\,XML};}

      \node[br,right=3mm] (@) at (@.east) {};
      \node (first-ast) at (@) {\usebox\FirstAst};
      \node[k] at (@.south) {Normalization};
      % \onslide<6->{\node[m] at (@.north) {in TypeScript};}
      \coordinate (@l) at(@.south west);

      \node[br,right=3mm] (@) at (@.east) {};
      \node (dataflow) at (@) {\usebox\DataFlow};
      \node[k] at (@.south) {Dataflow};

      \node[br,right=3mm] (@) at (@.east) {};
      \node (slicing) at (@) {\usebox\Slicing};
      \node[k] at (@.south) {Slicing};
      % \onslide<9->{\node[m] at (@.north) {Weiser~\cite{weiser_program_1984}};}

      \node[br,right=3mm] (@) at (@.east) {};
      \node (reconstruct) at (@) {\usebox\Reconstruct};
      \node[k] at (@.south) {Reconstruct};
      \coordinate (@r) at(@.south east);

      \scope[transparency group,a]
      \node[br,right=3mm] (@) at (@.east) {};
      \node (export) at (@) {\rotatebox{-5}{\usebox\ArchitectureBox}\llap{\T{.r\kern8pt}}};
      \node[k] at (@.south) {Output File};
      % \onslide<9->{\node[m] at (@.north) {Slice,~\ldots};}
      \endscope
   \draw[decorate,decoration={brace,mirror},line cap=round,line join=round] ([xshift=-1cm,yshift=-6mm]@l) to[edge node={coordinate (@flowr-location)}] ([xshift=1mm,yshift=-6mm]@r);
   \pgfonlayer{background}
      \node[below] (@flowr) at(@flowr-location) {\includegraphics[height=2\baselineskip]{data/flowR.pdf}};
      \fill[btdl@color@background,opacity=.5] (@flowr.north east) rectangle (@flowr.south west);
   \endpgfonlayer
\end{tikzpicture}}
\end{lrbox}

\begin{frame}[fragile,c]{}
\only<2->{\vspace*{-1em}}\only<-1|handout:0>{\vspace*{4em}}\begin{layout-full}\begin{center}%
\lstfs{9}
\begin{tikzpicture}
\node (@) at (0,0) {%
\solsetmintedstyle{plain number}%
\begin{minted}[deletekeywords={sum,prod},escapeinside={/*}{*/}]{R}
sum  <- 0
prod <- 1
n    <- 10

for(i in 1:(n-1)) {
   sum  <- sum + i
   prod <- prod * i
}

cat("Sum:", /*\T{\only<2->{\textbf}{sum}}*/, "\n")
cat("Product:", prod, "\n")
\end{minted}
};

   \node[right=3cm] (@2) at (@.east) {%
\solsetmintedstyle{plain}
\AnimateCode{onslide={o8:{7,5,4,3,1},-}, first slide=1,handout={1/1}}
\begin{minted}[deletekeywords={sum,prod},escapeinside={/*}{*/}]{R}
sum  <- 0
prod <- 1
n    <- 10

for(i in 1:(n-1)) {
   sum  <- sum + i
   prod <- prod * i
}

cat("Sum:", sum, "\n")
cat("Product:", prod, "\n")
\end{minted}
\endAnimateCode
   };
   \draw[-Kite,line cap=round] ([xshift=5mm]@.east) to[edge node={node[above] {\scriptsize\T{slice(\solGet{linenumbers}{\scriptsize\T{10}}, \textbf{sum})}}}] ([xshift=-5mm]@2.west);
\end{tikzpicture}
\begin{tikzpicture}[overlay,remember picture]
   \node[below=5mm,font=\sbfamily\Large] at(current page.north) {Program Slicing};
   \onslide<2->{
      \node[above] at(current page.south){\usebox\Overview};
   }
\end{tikzpicture}
\end{center}\end{layout-full}
\end{frame}

\begin{frame}
\begin{tikzpicture}[overlay,remember picture]
   % padding comes from preview border
   \node[above left=0mm] at (current page.south east) {%
      \includegraphics[width=.65\paperwidth]{build/visualizer-demo-graphic.pdf}
   };
   \node[below right=5mm,font=\sbfamily\Large] (@title) at(current page.north west) {%
      Interactive Visualization of Dataflow-Graphs%
   };
   \node[right=5mm,text width=.33\paperwidth] at(current page.west) {%
      \begin{itemize}
         \itemsep7.5pt
         \item 2\,--\,4 Students
         \item Visualize large graphs
         \item Web-Development
         \item Concepts:\\
            \textit{Dataflow-Graphs},\\\textit{Graph-Layouting},\\\textit{User Interfaces}
      \end{itemize}
   };
   \SupportedTypes
\end{tikzpicture}
\end{frame}

\begin{frame}
\begin{tikzpicture}[overlay,remember picture]
   \node[below right] at (current page.north west) {%
      \includegraphics[width=.65\paperwidth]{build/reconstruct-demo-graphic.pdf}
   };
   \node[above left=5mm,yshift=-1.5mm,font=\sbfamily\Large] (@title) at(current page.south east) {%
   Layout-Preserving Slice Reconstruction%
   };
   \node[above left,yshift=5mm,text width=.6\paperwidth] at(@title.north east) {%
      \begin{itemize}
         \item 1\,--\,3 Students
         \item Retain original layout
         \item Existing implementation in \href{https://www.typescriptlang.org/}{TypeScript}
         \item Concepts:\\
            \textit{Abstract Syntax Trees}, \textit{Parsers}, \textit{Pretty Printing}
      \end{itemize}
   };
   \SupportedTypes
\end{tikzpicture}
\end{frame}
\end{document}